\begin{enumerate}

\item Formar la matriz del tensor de Cauchy y obtener las tensiones principales en el plano xy.

\item Añadir la tensión normal sigma z, ordenar los tres autovalores y calcular la máxima cortadura absoluta.

\item Obtener la tensión equivalente de von Mises y la magnitud equivalente de Tresca.

\item Aplicar las relaciones de transformación al elemento girado.

\item Utilizar la ley de Hooke generalizada para deformaciones y distorsión.

\end{enumerate}

\noindent\textbf{Operaciones y resultados intermedios.}

\begin{enumerate}

\item Centro del círculo en el plano xy.

\[\fitmath{y_{1}=\frac{110+45}{2}\approx 77.5\,\mathrm{MPa}}\]

\item Radio del círculo de Mohr.

\[\fitmath{y_{2}=\sqrt{{\left(\frac{110-45}{2}\right)}^{2}+{35}^{2}}\approx 47.76243\,\mathrm{MPa}}\]

\item Tensión principal mayor en el plano.

\[\fitmath{y_{3}=\frac{110+45}{2}+\sqrt{{\left(\frac{110-45}{2}\right)}^{2}+{35}^{2}}\approx 125.2624\,\mathrm{MPa}}\]

\item Tensión principal menor en el plano.

\[\fitmath{y_{4}=\frac{110+45}{2}-\sqrt{{\left(\frac{110-45}{2}\right)}^{2}+{35}^{2}}\approx 29.73757\,\mathrm{MPa}}\]

\item Von Mises a partir de las componentes.

\[\fitmath{y_{5}=\frac{\sqrt{{\left(110-45\right)}^{2}+{\left(45-0\right)}^{2}+{\left(0-110\right)}^{2}+6\cdot {35}^{2}}}{\sqrt{2}}\approx 113.3578\,\mathrm{MPa}}\]

\item Índice tensional de Tresca.

\[\fitmath{y_{6}=2\cdot \left(\frac{5\cdot \sqrt{365}}{4}+\frac{155}{4}\right)\approx 125.2624\,\mathrm{MPa}}\]

\item Tensión normal sobre la cara x prima.

\[\fitmath{y_{7}=\frac{110+45}{2}+\frac{110-45}{2}\cdot \cos\left(\frac{2\cdot \pi\cdot 45}{180}\right)+35\cdot \sin\left(\frac{2\cdot \pi\cdot 45}{180}\right)\approx 112.5\,\mathrm{MPa}}\]

\item Tensión normal sobre la cara y prima.

\[\fitmath{y_{8}=\frac{110+45}{2}-\frac{110-45}{2}\cdot \cos\left(\frac{2\cdot \pi\cdot 45}{180}\right)-35\cdot \sin\left(\frac{2\cdot \pi\cdot 45}{180}\right)\approx 42.5\,\mathrm{MPa}}\]

\item Cortante sobre el plano girado.

\[\fitmath{y_{9}=-\left(\frac{110-45}{2}\right)\cdot \sin\left(\frac{2\cdot \pi\cdot 45}{180}\right)+35\cdot \cos\left(\frac{2\cdot \pi\cdot 45}{180}\right)\approx -32.5\,\mathrm{MPa}}\]

\item Deformación epsilon x.

\[\fitmath{y_{10}=\frac{110-0.33\cdot \left(45+0\right)}{71700}\approx 0.001327057\,}\]

\item Deformación epsilon y.

\[\fitmath{y_{11}=\frac{45-0.33\cdot \left(110+0\right)}{71700}\approx 0.0001213389\,}\]

\item Deformación epsilon z.

\[\fitmath{y_{12}=\frac{0-0.33\cdot \left(110+45\right)}{71700}\approx -0.0007133891\,}\]

\item Distorsión ingenieril gamma xy.

\[\fitmath{y_{13}=\frac{2\cdot \left(1+0.33\right)\cdot 35}{71700}\approx 0.001298466\,}\]

\item Respuesta principal: comprobación final.

\[\fitmath{y_{14}=\frac{0-0.33\cdot \left(110+45\right)}{71700}\approx -0.0007133891\,}\]

\end{enumerate}
